Edge caching reduces delay and backhaul load by placing content near users, but IoT edge networks are difficult caching environments because user mobility changes both demand and the effective cooperation topology among edge nodes. We study topology-aware cooperative caching in a graph-based IoT edge system where access points act as edge caches and mobility-induced handoffs define weighted inter-node relationships. We formulate cache placement as a stochastic online optimization problem that jointly considers access cost, replication overhead, topology-aware redundancy, relocation cost, and optimization overhead under random topology perturbations. To solve the problem, we propose Adaptive TACC, a reproducible pipeline that combines topology-aware greedy placement, transition-aware DQN refinement, and a cost-aware validation selector that changes the deployed placement only when the expected benefit justifies cache rewriting. Because the Dartmouth campus WLAN trace is publicly described but not directly downloadable through legacy unauthenticated URLs at the time of this study, we record the official dataset access path and use a deterministic trace-compatible mobility generator for reproducible evaluation. Experiments over a 20-node mobility graph, 48-object catalog, and repeated perturbation samples show that DQN is useful when it is used as an online local-refinement operator rather than as a fixed alternative placement. Adaptive TACC selects online DQN refinement in all four tested perturbation windows, improving objective from 2.633 to 2.209 in the stable window and from 5.252 to 3.555 at \(p=0.15\), while explicitly reporting the relocation cost of each cache update.
